La tabella seguente riassume la provvigione attuale di ciascuna compresa;
seguono poi le stime suddivise per singole particelle. La metodologia di calcolo
è quella descritta nella sezione~\ref{sec:note-volumi}, basata su
Tabacchi~\cite{tabacchi:2011}. La provvigione comprende solo le piante mature
(classe diametrica 20 cm e superiori), le stesse considerate per il calcolo
della ripresa; le piante sottomisura sono escluse. Le colonne ``IF inf'' e ``IF
sup'' denotano i limiti inferiori e superiori, rispettivamente, dell'intervallo
fiduciario al 95\%.

@@provvigione(alberi=alberi-calcolati.csv,per_compresa=si,per_particella=no,per_genere=no,intervallo_fiduciario=si,totali=si)

\subsection*{Provvigione: Serra}
@@provvigione(alberi=alberi-calcolati.csv,compresa=Serra,per_compresa=no,per_particella=si,per_genere=no,intervallo_fiduciario=si,totali=si)
\clearpage

\subsection*{Provvigione: Fabrizia}
@@provvigione(alberi=alberi-calcolati.csv,compresa=Fabrizia,per_compresa=no,per_particella=si,per_genere=no,intervallo_fiduciario=si,totali=si)

\subsection*{Provvigione: Capistrano}
@@provvigione(alberi=alberi-calcolati.csv,compresa=Capistrano,per_compresa=no,per_particella=si,per_genere=no,intervallo_fiduciario=si,totali=si)

\clearpage
